% part: intuitionistic-logic
% chapter: propositions-as-types
% section: introduction

\documentclass[../../../include/open-logic-section]{subfiles}

\begin{document}

\olfileid[pt]{int}{pty}{int}

\olsection{Introdução}

O cálculo lambda é um modelo simples de computação. Ele opera sobre
termos construídos a partir de variáveis. Se $N$ e $M$ são termos,
então $(NM)$ também é um termo (``$N$ aplicado a~$M$''). Se $N$ é um
termo, então $\lambd[x][N]$ é um termo. Se $N$ contém a variável~$x$,
então as ocorrências correspondentes de~$x$ em $\lambd[x][N]$ são
ligadas pelo operador $\lambd[x]$ e, portanto, deixam de ser livres.
Dizemos que um termo da forma $(\lambd[x][N]M)$ se $\beta$-contrai a
$\Subst{N}{M}{x}$ (o resultado de substituir toda ocorrência livre de
$x$ em~$N$ por $M$). Um termo $N$ se $\beta$-converte em~$N'$ se $N'$
resulta de $N$ pela $\beta$-contração de um subtermo, e escrevemos $N
\redone N'$. Uma computação no $\lambd$-cálculo é uma sequência $N
\redone N_1 \redone N_2 \redone \dots$. Se essa sequência termina em
um termo $N_k$ no qual nenhum subtermo pode ser $\beta$-contraído, ele
é chamado de \emph{normal}, ou de forma normal de~$N$. Pense nas
computações no $\lambd$-cálculo como tais sequências. A computação
termina se resultar em uma forma normal, e a forma normal é o resultado da
computação. Esse modelo simples de computações revela-se completo no
sentido de Turing: toda função recursiva parcial pode ser computada
por um termo no $\lambd$-cálculo.

Haskell Curry e William Howard descobriram, independentemente, uma
estreita semelhança entre a dedução natural e o $\lambd$-cálculo.
Intuitivamente, um termo $\lambd[x][N]$ é uma função: $x$ é o argumento;
$N$ nos diz como calcular o valor dado~$x$. O termo
$(\lambd[x][N]M)$ consiste então em aplicar a função $\lambd[x][N]$ a
um argumento~$M$, e a $\beta$-contração nos diz como avaliá-lo:
substitua $x$ em $N$ por~$M$ (e continue avaliando o termo resultante).
Agora pense nessas expressões como funções com domínio e imagem fixos.
Se o domínio de $\lambd[x][N]$ é $A$ e e sua imagem é~$B$, então
devemos pensar que a variável~$x$ recebe apenas argumentos de~$A$ e
que $N$ produz valores de~$B$. Consequentemente, $(\lambd[x][N]M)$ só
deve fazer sentido se $\lambd[x][N]$ é uma função $A \to B$ e $M \in
A$. Se quisermos acrescentar essa informação ao $\lambd$-cálculo,
obteremos a versão \emph{tipada} do $\lambd$-cálculo. No
$\lambd$-cálculo tipado,
nem toda expressão $(\lambd[x][N]M)$ é permitida, mas apenas aquelas
em que os ``tipos'' de $\lambd[x][N]$ e $M$ se ajustam, isto é,
$\lambd[x][N]$ tem tipo $A \lif B$ e $M$ tem tipo~$A$. O tipo de
$\lambd[x][N]$ é $A \to B$ somente se o tipo de $x$ é $A$ e o tipo de
$N$ é~$B$. Em geral, nem toda expressão $(M_1M_2)$ é permitida. São
permitidos apenas os termos $(M_1M_2)$ em que $M_1$ tem tipo $A \to B$
e $M_2$ tem tipo~$A$; nesse caso, $(M_1M_2)$ tem tipo~$B$.

Essas regras de tipagem correspondem agora claramente a
!!{derivation}s em dedução natural, nas quais os tipos são
representados por !!{formula}s e as próprias !!{derivation}s por
$\lambd$-termos. Por exemplo, !!a{derivation} de $!A \lif !B$
corresponde a um termo $\lambd[\typeof{x}{!A}][N]$, que tem o tipo de
função $!A \lif !B$. As regras de inferência da dedução natural
correspondem a regras de tipagem no cálculo lambda tipado, por exemplo,
\begin{prooftree}
  \AxiomC{$\Discharge{!A}{x}$}
  \DeduceC{$!B$}
  \DischargeRule{\Intro\lif}{x}
  \UnaryInfC{$!A \lif !B$}
  \DisplayProof\bottomAlignProof
  \quad\text{corresponde a}\quad
  \Axiom$x: !A \fCenter N: !B$
  \UnaryInf$ \fCenter \lambd[\typeof{x}{!A}][N] : !A \lif !B$
\end{prooftree}
em que a regra à direita significa que, se $x$ é do tipo~$!A$ e $N$ é
do tipo~$!B$, então $\lambd[\typeof{x}{!A}][N]$ é do tipo~$!A \lif !B$.

A regra $\Elim\lif$ corresponde à regra de tipagem para termos de
composição, isto é,
\[
  \AxiomC{$!A \lif !B$}
  \AxiomC{$!A$}
  \RightLabel{\Elim\lif}
  \BinaryInfC{$!B$}
  \DisplayProof
  \quad\text{corresponde a}\quad
  \Axiom$\fCenter M_1: !A \lif !B$
  \Axiom$\fCenter M_2: !A$
  \BinaryInf$ \fCenter (M_1M_2) : !B$
  \DisplayProof
\]
Se uma regra \Intro\lif{} é seguida imediatamente por uma regra
\Elim\lif{}, a !!{derivation} pode ser simplificada:
\begin{gather*}
  \AxiomC{$\Discharge{!A}{x}$}
  \DeduceC{$!B$}
  \DischargeRule{\Intro\lif}{x}
  \UnaryInfC{$!A \lif !B$}
  \AxiomC{}
  \DeduceC{$!A$}
  \RightLabel{\Elim\lif}
  \BinaryInfC{$!B$}
  \DisplayProof
  \quad\redone\quad
  \AxiomC{}
  \DeduceC{$!A$}
  \DeduceC{$!B$}
  \DisplayProof
\end{gather*}
o que corresponde à $\beta$-contração de termos lambda
\[
(\lambd[\typeof{x}{!A}][N]M) \quad\redone\quad
\Subst{N}{M}{x}.
\]

Correspondências semelhantes valem entre as regras para $\land$ e os
tipos ``produto'', e entre as regras para $\lor$ e os tipos ``soma''.

Essa correspondência entre termos no cálculo lambda simplesmente
tipado e !!{derivation}s em dedução natural é chamada de
correspondência ``Curry--Howard'', ``provas como programas'' ou
``proposições como tipos''. Além das !!{formula}s (proposições)
corresponderem a tipos, e das provas corresponderem a termos, podemos
resumir as correspondências da seguinte maneira:
\begin{center}
  \begin{tabular}{c c c}
    lógica & programa \\
    \hline
    proposição/!!{formula} & tipo \\
    prova & termo \\
    hipótese & variável \\
    hipótese descarregada & variável ligada \\
    hipótese não descarregada & variável livre \\
    $!A \lif !B$ & tipo de função \\
    $!A \land !B$ & tipo produto \\
    $!A \lor !B$ & tipo soma \\
    $\lfalse$ & tipo vazio \\
    \hline
  \end{tabular}
\end{center}

A correspondência de Curry--Howard é um dos pilares dos assistentes
automatizados de demonstração e dos verificadores de tipos para
programas, pois verificar uma prova que testemunha uma proposição
(como fizemos acima) equivale a verificar se um programa (termo) tem o
tipo declarado.
\end{document}
